\section{Machine Learning in Sports Analytics}
\label{sec:ml-sports}

\subsection{From Regression to Ensemble Methods}

The application of statistical modeling to sports outcomes has a history spanning several decades, but the modern era of machine-learning-driven sports analytics is considerably more recent. Early work in sports prediction relied on traditional regression techniques: logistic regression for binary outcomes (win/loss), linear regression for continuous metrics (points, yards), and occasionally Poisson regression for count data (goals scored) \parencite{boulier2003predicting}. These parametric approaches offered interpretability and statistical inference but imposed strong functional-form assumptions that limited their ability to capture the complex, non-linear relationships inherent in athletic competition.

The transition from parametric regression to machine learning in sports analytics accelerated in the 2010s as three developments converged. Granular play-by-play data  -- previously locked behind proprietary vendor agreements or manual charting  -- became widely available as leagues invested in tracking infrastructure and open-source data pipelines emerged \parencite{carl2021nflfastr}. Simultaneously, the computational cost of training ensemble models dropped by orders of magnitude, making random forests and gradient boosting accessible to individual researchers without institutional computing clusters. And the success of ensemble methods in structured prediction competitions (notably Kaggle) demonstrated their superiority over parametric alternatives for tabular prediction tasks with heterogeneous features and non-linear relationships \parencite{chen2016xgboost}.

\textcite{friedman2001gradientboosting} introduced the gradient boosting framework that would become foundational to modern sports analytics. By framing ensemble construction as iterative gradient descent in function space, Friedman unified previously disparate boosting approaches under a single theoretical umbrella and demonstrated that the resulting models could achieve state-of-the-art performance on a wide variety of supervised learning tasks. The framework's key insight  -- that weak learners (shallow trees) can be combined additively to approximate arbitrarily complex target functions  -- proved particularly well-suited to sports data, where outcome probabilities depend on high-order interactions between situational variables that cannot be specified a priori.

\textcite{chen2016xgboost} extended Friedman's framework with computational optimizations (column subsampling, cache-aware parallel tree construction) and regularization innovations (L1 and L2 penalties on leaf weights, maximum depth constraints) that made gradient boosting practical for large-scale applications while reducing overfitting. XGBoost became the dominant algorithm for tabular prediction tasks and was rapidly adopted by the sports analytics community for problems ranging from player valuation to game outcome prediction. Its combination of predictive power, native handling of heterogeneous features, and built-in regularization made it particularly attractive for sports applications where sample sizes are modest by machine learning standards (an NFL season produces only ${\sim}270$ games) and overfitting is a constant concern.

\subsection{Expected Goals in Association Football}
\label{sec:xg-literature}

The most direct methodological antecedent to the \xscore{} model is the Expected Goals (xG) framework in association football, which emerged as a practitioner tool in the early 2010s  -- circulating on analytics blogs and within club data departments from approximately 2012 onward  -- and has since become the standard metric for evaluating team and player performance in European football leagues \parencite{eggels2016expected}. The xG approach assigns to each shot a probability of resulting in a goal, conditional on situational features describing the shot's location, angle, body part used, preceding play pattern, and defensive pressure. By summing these probabilities across all shots in a match, analysts obtain an expected goal total for each team that represents what a league-average finishing team would have scored given the same sequence of chances. The difference between actual goals and expected goals then isolates finishing quality from chance creation, providing a decomposition of offensive output into its constituent parts.

\textcite{rathke2017expectedgoals} provided one of the early peer-reviewed \emph{academic} examinations of xG methodology, analyzing shot efficiency in professional football using logistic regression to estimate goal probability from shot location and type. While Rathke's model was methodologically simple by contemporary standards  -- using only distance from goal, shot angle, and a header indicator as features  -- it established the conceptual framework that would be refined by subsequent work: a calibrated probability baseline enables meaningful performance measurement by separating what happened from what was expected to happen given the situation. This ``actual minus expected'' logic is precisely the architecture underlying the Game Deserved Score framework developed in this paper, transposed from the shot level in football to the drive level in American football.

Subsequent developments in xG modeling moved beyond logistic regression to ensemble methods. Industry models developed by analytics companies such as StatsBomb \parencite{statsbomb2019opendata}, Opta, and Understat employ gradient boosted trees or neural networks trained on datasets comprising millions of shots, incorporating features for defensive positioning, game state, shot sequence patterns, and player-specific finishing ability. The academic literature has largely validated these approaches: studies consistently find that xG-based metrics outperform raw goal counts for predicting future team performance \parencite{rathke2017expectedgoals,anzer2021goal}, that the expected-actual gap captures meaningful team quality differences \parencite{eggels2016expected}, and that calibration quality directly affects the utility of the resulting performance metrics.

The xG framework illustrates several design principles that I adopt in the present work. The probability baseline must exclude team identity from its features  -- an xG model that knows which player is shooting will produce higher expected goals for elite finishers, compressing the actual-minus-expected gap toward zero and obscuring the quality signal. The unit of analysis must balance signal stability against contextual granularity: shots are the natural unit for association football because each shot is an independent scoring opportunity with well-defined situational features, just as drives serve this role in American football. And calibration is not merely desirable but essential, because the downstream metric treats predicted probabilities as genuine expectations from which deviations carry quantitative meaning.

\subsection{Machine Learning in American Football}
\label{sec:ml-nfl}

The application of machine learning to NFL outcome prediction has followed a distinct trajectory from the xG revolution in European football, characterized by later adoption, greater methodological fragmentation, and a sharper divide between public and proprietary analytics. While Major League Baseball experienced its sabermetric revolution in the early 2000s \parencite{lewis2003moneyball} and association football's xG framework matured by 2015, the NFL analytics community did not coalesce around machine-learning-based frameworks until the late 2010s. Earlier statistical approaches to NFL prediction relied on linear models and Bayesian methods \parencite{harville1980predictions,glickman1998statespacenfl}, which captured opponent quality and temporal dynamics but lacked the flexibility to model the high-dimensional interaction structures inherent in play-by-play data.

\textcite{lock2014randomforests} represents one of the earliest peer-reviewed applications of ensemble machine learning to NFL game-level prediction. Their work employed random forests \parencite{breiman2001randomforests} to estimate win probability before each play, demonstrating that tree-based methods could capture non-linear interactions between game-state variables (score differential, time remaining, field position, possession) that eluded parametric win-probability models \parencite{stern1991probability,glickman1998statespacenfl}. Critically, Lock and Nettleton found that offensive variables  -- particularly passing efficiency metrics  -- contributed more to win-probability predictions than defensive variables, an observation consistent with the offensive primacy finding of this paper, though they did not frame their results as a test of the ``defense wins championships'' hypothesis.

\textcite{yurko2019nflwar} developed the nflWAR framework, which adapted the Wins Above Replacement (WAR) concept from baseball to American football using a multi-level model built on play-by-play EPA data. Their contribution was methodological rather than substantive: they demonstrated that individual player contributions could be isolated from team effects using hierarchical Bayesian modeling applied to NFL play-by-play data. The nflWAR framework established the precedent  -- followed by this paper  -- that open-source, reproducible analytical tools built on publicly available data can produce insights competitive with proprietary systems. However, nflWAR operated at the player level rather than the team level, and its use of EPA as the underlying currency inherits the play-level noise and context-sensitivity issues discussed in \Cref{sec:epa-framework}.

The \texttt{nflfastR} package \parencite{carl2021nflfastr} transformed the landscape of public NFL analytics by providing a clean, well-documented, and continuously maintained R interface to official NFL play-by-play data. By democratizing access to the raw data underlying EPA calculations, \texttt{nflfastR} enabled a proliferation of independent analyses that had previously been impossible without proprietary data access. The present paper builds directly on this foundation: all 46{,}920 drives in the analytical sample (training and test combined) are constructed from \texttt{nflfastR} play-by-play records, and the model's reproducibility guarantee rests on the public availability of this data source.

\subsection{Probability Calibration in Predictive Sports Models}
\label{sec:calibration-literature}

A recurring theme across the machine learning sports analytics literature is the distinction between discrimination (ranking observations correctly) and calibration (producing accurate probability estimates). \textcite{niculescumizil2005probabilities} demonstrated that many supervised learning algorithms  -- including support vector machines, random forests, and boosted decision trees  -- achieve strong discrimination while producing poorly calibrated probability outputs. Their finding that boosted tree models tend to push predictions toward 0 and 1 (over-confidence near probability boundaries) motivated the development of post-hoc calibration techniques, including Platt scaling and isotonic regression.

This finding has particular relevance for sports analytics applications where the predicted probability is not merely used for ranking but serves as a quantitative baseline against which performance is measured. In an xG model, a miscalibrated probability of 0.40 assigned to a shot that should receive 0.25 will systematically undercount the finishing quality of teams that score on such chances and overcredit teams that miss them. The distortion propagates through every downstream calculation that treats the prediction as an expectation. The same logic applies with even greater force to the \xscore{} model, where calibrated drive-level probabilities serve as the denominator in the \gds{} framework's expected-value calculations. \textcite{niculescumizil2005probabilities} found that isotonic regression  -- a non-parametric calibration method that fits a monotone step function from raw scores to calibrated probabilities  -- consistently outperformed Platt scaling for tree-based models, a finding that directly motivates the calibration approach adopted in \Cref{sec:calibration}.

\textcite{cervone2016multiresolution} demonstrated the value of calibrated probability models in basketball analytics, developing a multiresolution stochastic process model for predicting possession outcomes in the NBA. Their approach decomposed each possession into a sequence of decision points, assigning transition probabilities at each stage, and aggregated these into a possession-level expected value. While operating in a different sport and at a different temporal resolution, their work shares a core methodological commitment with the present paper: that calibrated probability estimates at a natural unit of play (possessions in basketball, drives in football) provide a more stable and interpretable basis for performance measurement than play-by-play aggregation. The success of their approach in capturing genuine team quality differences reinforced the viability of possession-level (or drive-level) expected-value frameworks.
